\documentclass[11pt,a4paper]{article}
\usepackage[utf8]{inputenc}
\usepackage[margin=1in]{geometry}
\usepackage{amsmath}
\usepackage{amsfonts}
\usepackage{amssymb}
\usepackage{hyperref}
\usepackage{booktabs}
\usepackage{graphicx}
\usepackage{listings}
\usepackage{xcolor}

% ── Mermaid listing style ──────────────────────────────────────────
\definecolor{mermaidbg}{HTML}{F9F6F1}
\lstdefinelanguage{mermaid}{
  morekeywords={graph, subgraph, end, TD, direction, TB},
  sensitive=true,
  morecomment=[l]{\%\%},
  morestring=[b]{"},
}
\lstdefinestyle{mermaidstyle}{
  language=mermaid,
  backgroundcolor=\color{mermaidbg},
  basicstyle=\ttfamily\footnotesize,
  breaklines=true,
  frame=single,
  rulecolor=\color{gray},
  captionpos=b,
}

\title{\textbf{Optimized Topic Modeling of \\AI Research Literature:
Tracking Thematic Evolution and Trends in Nepal }}
\author{
  Bibek Paudyal \quad (Softwarica ID: 250288) \\
  Siddhartha Bhatta \quad (Softwarica ID: 250620) \\
  Sajan Mahat \quad (Softwarica ID: 250289) \\[1em]
  \text{Module: ST7085CEM - Advanced Machine Learning} \\
  \text{Softwarica College of IT \& E-Commerce}\\
  \text{(Affiliated to Coventry University)}
}
\date{\today}

\begin{document}

% ── Title spans both columns, then switch to two-column ────────────
\twocolumn[
  \maketitle
]

% ====================================================================
\section{Research Title / Topic}
% ====================================================================
Optimized Topic Modeling of AI Research Involving Nepali Authors:
Tracking Thematic Evolution and Trends in Nepal (2015--2025)

% ====================================================================
\section{Introduction \& Background}
% ====================================================================
Artificial Intelligence (AI) research has grown exponentially worldwide,
and Nepal's academic and research community has actively contributed to
this global landscape through localized applications, natural language
processing (NLP) for low-resource languages, and computer vision
advancements. While individual studies document specific subfields,
tracking the macro-level evolution of AI research produced by Nepali
authors over the past decade remains a challenge. Traditional literature
reviews are slow, manually intensive, and prone to individual bias.

To overcome this, unsupervised topic modeling provides a powerful tool to
discover hidden themes. While classical probabilistic models like Latent
Dirichlet Allocation (LDA) offer high interpretability, they suffer from
hyperparameter sensitivity and instability. Conversely, modern
embedding-based models like BERTopic capture semantic relationships
effectively. This study applies an evolutionary-optimized topic modeling
framework specifically to peer-reviewed literature authored or
co-authored by researchers affiliated with Nepal, providing a
comprehensive view of how AI research has evolved locally from 2015 to
2025.

% ====================================================================
\section{Problem Statement}
% ====================================================================
Although global AI literature is heavily analyzed, there is a lack of
systematic, automated scientometric evaluation tracking how AI research
involving Nepali authors and researchers has evolved over time.
Furthermore, conventional topic modeling relies on unoptimized classical
parameters, and there is a need to compare whether
evolutionary-optimized LDA can outperform or complement modern neural
methods (such as BERTopic) when applied to regional academic corpora.

% ====================================================================
\section{Research Objective}
% ====================================================================
Compare classical (LDA, GA-optimized LDA) and modern (BERTopic) topic
modeling on one time-stamped corpus of AI research, discovering themes,
quantifying their evolution, and judging which approach models the field
best.

% ====================================================================
\section{Brief Literature Review}
% ====================================================================
Recent literature highlights a growing shift toward optimizing topic
discovery. Agrawal et al.\ (2018) demonstrated that metaheuristic
optimization fixes LDA instability. Yu \& Xiang (2023) applied standard
LDA to global AI research literature, establishing a baseline for
thematic mapping. Furthermore, recent studies (e.g., 2025 comparative
evaluations of neural and classical topic models) indicate mixed
performance results between embedding-based techniques and probabilistic
baselines across various domains, justifying a comprehensive comparative
benchmark. Specifically focusing this approach on regional contributions
(such as Nepal-affiliated research) fills a distinct gap in localized
scientometric tracking.

% ====================================================================
\section{Proposed Methodology}
% ====================================================================
The research follows a six-stage iterative computational pipeline, as
illustrated in Figure~\ref{fig:pipeline}.

\begin{enumerate}
    \item \textbf{Corpus Construction.}
          Build a time-stamped, peer-reviewed corpus of AI research
          papers (titles and abstracts) spanning 2015 to 2025 with at
          least one author affiliated with a Nepali institution.
          Papers are retrieved from the Semantic Scholar API with
          supplementary harvesting from NepJOL
          (\url{https://www.nepjol.info/}), IEEE Xplore, and Springer,
          then de-duplicated and venue-verified.

    \item \textbf{Preprocessing.}
          Implement domain-aware text preprocessing, tokenization,
          lemmatization, stop-word removal, and Vector Space Model
          (VSM) Bag-of-Words / TF-IDF generation, while preserving
          critical technical acronyms and localized terms.

    \item \textbf{Model Implementation.}
          Fit three models onto the identical Nepal-focused corpus in
          parallel:
          \begin{itemize}
              \item \textit{Standard LDA} ,  probabilistic baseline
                    with manual or grid-search settings.
              \item \textit{GA-Optimized LDA} ,  a Genetic Algorithm
                    iteratively optimizes hyperparameters
                    ($K, \alpha, \beta$) by maximizing a multi-objective
                    fitness function combining topic coherence ($C_v$),
                    diversity, and stability.
              \item \textit{BERTopic} ,  sentence-transformer
                    embeddings + HDBSCAN clustering to capture semantic
                    nuances in local research.
          \end{itemize}

    \item \textbf{Multi-Dimensional Evaluation.}
          Score all models rigorously using semantic coherence ($C_v$),
          topic diversity, and stability metrics to determine which
          approach models the corpus best.

    \item \textbf{Temporal Evolution Tracking.}
          Slice the corpus into temporal windows
          (2015--17, 2018--20, 2021--23, 2024--25) and re-fit the
          best-performing model(s) per window to map how local AI
          research themes shift across the decade.

    \item \textbf{Automated Labeling.}
          Leverage Large Language Model (LLM) assistance to generate
          coherent, human-interpretable thematic labels for every
          discovered topic.
\end{enumerate}

% ── Pipeline diagram (Mermaid source) ──────────────────────────────
% Render the Mermaid code below with the Mermaid CLI or online editor
% (https://mermaid.live) and export as pipeline.png, then uncomment
% the \includegraphics line.

\begin{figure*}[htbp]
  \centering
  \includegraphics[width=0.88\textwidth]{ga_lda_architecture.png}
  \caption{End-to-end methodology pipeline.}
  \label{fig:pipeline}
\end{figure*}


% ====================================================================
\section{Dataset Description}
% ====================================================================
The corpus comprises peer-reviewed AI/ML research papers (titles and
abstracts) spanning from 2015 to 2025. Papers are collected from
multiple sources:

\begin{itemize}
    \item \textbf{Semantic Scholar API}
          (\url{https://api.semanticscholar.org/})  -  primary
          retrieval and venue-verification engine.
    \item \textbf{NepJOL}
          (\url{https://www.nepjol.info/})  -  Nepal Journals Online,
          aggregating Nepali academic publications across disciplines.
    \item \textbf{IEEE Xplore}
          (\url{https://ieeexplore.ieee.org/})  -  conference and
          journal papers in electrical engineering, computer science,
          and AI. Nepal-affiliated papers are retrieved via the
          Advanced Search command
          \texttt{"Author Affiliations":Nepal}
          (\url{https://ieeexplore.ieee.org/search/advanced}).
    \item \textbf{Springer}
          (\url{https://link.springer.com/})  -  internationally
          indexed research with Nepali-affiliated authorship,
          filtered by searching with Nepal-based institution names
          in the author affiliation facet.
\end{itemize}

\noindent
Inclusion is strictly filtered to papers containing $\ge 1$ author with
a Nepal-based institutional affiliation (e.g., Tribhuvan University,
Kathmandu University, NAAMII, etc.), ensuring the dataset directly
reflects domestic and diaspora contributions to AI from Nepal.


% ====================================================================
\section{Proposed Machine Learning Model}
% ====================================================================
\begin{itemize}
    \item \textbf{Baseline 1:} Standard LDA (probabilistic baseline
          with manual or grid settings).
    \item \textbf{Proposed Model:} GA-Optimized LDA (metaheuristic
          multi-objective search tuning topic count and Dirichlet
          priors).
    \item \textbf{Benchmark Model:} BERTopic (utilizing sentence
          transformers for dense vector spaces and HDBSCAN clustering
          to capture semantic nuances in local research).
\end{itemize}

% ====================================================================
\section{Expected Results}
% ====================================================================
\begin{itemize}
    \item A comparative evaluation table measuring LDA, GA-LDA, and
          BERTopic across coherence, diversity, and interpretability
          for the Nepal-focused dataset.
    \item Longitudinal trend visualizations detailing how AI research
          themes in Nepal shifted between 2015 and 2025 (e.g.,
          tracking the transition from traditional machine
          learning/local NLP tasks toward deep learning, transformer
          models, and generative AI applications).
    \item A validated, labeled topic set providing clear thematic
          breakdowns of the trajectory of AI research in Nepal.
\end{itemize}

% ====================================================================
\section{References}
% ====================================================================
\begin{enumerate}
    \item Agrawal, A., Fu, W., \& Menzies, T. (2018). What is Wrong
          with Topic Modeling? (and How to Fix it Using Search-based
          Software Engineering). \textit{Information and Software
          Technology}.
    \item Yu, J., \& Xiang, Y. (2023). Discovering Topics and Trends
          in the Field of Artificial Intelligence: Using LDA Topic
          Modeling. \textit{Expert Systems with Applications}.
    \item Mapping the Evolution of Artificial Intelligence in
          Agriculture: A Large-Scale BERTopic Analysis.
          \textit{MDPI AgriEngineering}, 8(8), 312.
    \item Haq, M.~I.~U., Li, Q., \& Hassan, S. (2019). Text Mining
          Techniques to Capture Facts for Cloud Computing Adoption
          and Big Data Processing. \textit{IEEE Access}, 7,
          162254--162267.
\end{enumerate}


\end{document}
